\documentclass[11pt,a4paper]{article}
\usepackage[utf8]{inputenc}
\usepackage[T1]{fontenc}
\usepackage{geometry}
\geometry{left=2.7cm,right=2.7cm,top=2.6cm,bottom=2.6cm}
\usepackage{lmodern,microtype,setspace}
\usepackage{graphicx,float,amsmath,amssymb,cite,listings,xcolor,booktabs,array,caption,subcaption,algorithm,algpseudocode,tikz,pgfplots,seqsplit}
\usepackage[hidelinks]{hyperref}
\pgfplotsset{compat=1.18}
\usetikzlibrary{shapes,arrows,positioning,fit,backgrounds,matrix,calc}
\onehalfspacing
\setlength{\parindent}{1.5em}
\setlength{\parskip}{0.15em}
\hypersetup{
  pdftitle={Privacy Lens: An Evidence-Bounded Android Privacy Review Framework, Revision 25},
  pdfauthor={Zhiheng Zhang}
}
\title{Privacy Lens: An Evidence-Bounded Android Privacy Review Framework with Independently Witnessed Trust-Store Releases (Revision 25)}
\author{Zhiheng Zhang \\ Student ID: 54560484 \\ Supervisor: Richard Sinnott}
\date{August 2026}
\begin{document}
\maketitle
\begin{abstract}
\small\singlespacing
Mobile privacy-review tools face an evidence-to-decision gap: a permission observation can indicate activity but cannot by itself establish purpose, necessity, lawful basis, acquisition completeness, or legal compliance. This thesis addresses that gap through Privacy Lens, a local-first Android prototype that converts bounded permission-audit summaries into traceable prompts for human review.

Following a design-science method, the architecture treats provenance, temporal scope, missing context, and uncertainty as first-class evidence rather than converting a count into a legal verdict. A fail-closed engine rejects malformed records and unsafe temporal aggregation, isolates simulator data, and records the rule-pack version applied to each finding. Jurisdiction mappings remain outside the Android, repository, and interface layers through a replaceable pack contract; the European Union General Data Protection Regulation provides the evaluated legal objective.

Revision 25 prevents a root-signed reviewer trust-store envelope from becoming current by signature alone. After the existing canonical Ed25519, validity, sequence, predecessor, revocation, and persisted rollback checks, a separately configured policy requires at least two distinct witness keys and enough valid receipts to meet its threshold. Each receipt binds the exact envelope digest, sequence, witness date, and declared application release identity. Missing or invalid policy, receipts, or quorum exposes no reviewer anchor and advances no rollback state. Production roots, envelopes, witness policy, receipts, and legal attestation remain empty, so the evaluated EU pack remains fail-closed.

The fixed-seed 1,000-case corpus returned 800 true positives and 200 true negatives, while a further 1,800 independently calculated governance cases and the expanded witness adversarial probes passed. Accessibility, release-privacy, compilation, lint, APK/AAB build, local APK signature verification, and a short physical-device sample also passed within their stated scopes. App 1.13.0 showed witnesses 0/2, an unprovisioned production policy, and the limits of artifact identity, witness identity, independence, legal correctness, and transparency without observed clipping; three cold launches returned a live process and the crash buffer remained empty.

These results support a stronger engineering and publication candidate, internally assessed at 96/100 against the recorded 70/100 predecessor baseline. They do not establish qualified legal review, production witness governance, installed-artifact identity, public consistency, unrestricted Android visibility, accessibility conformance, participant outcomes, long-run reliability, store approval, or publication acceptance. NIST, TUF, and in-toto concepts inform bounded design choices, but Privacy Lens claims neither derived GDPR obligations nor conformance to those systems. The contribution is an evidence-bounded architecture that makes both cryptographic progress and residual authority limits visible.
\end{abstract}
\newpage
\tableofcontents
\newpage
\input{Iteration-V2/Chapter1/chapter-01-23}

\subsection{Revision 24 contribution: a signed monotonic trust-store boundary}

Revision 24 replaces the preceding static reviewer-key denylist concept with a canonical, Ed25519-signed trust-store envelope. The envelope identifies an offline signing root, separates issuance from approval, carries reviewer anchors and explicit revocations, expires, advances through a positive safe-integer sequence, and binds every non-initial version to the digest of its accepted predecessor. The application persists the highest accepted sequence and envelope digest and refuses older, forked, skipped, or unchained candidates. Production roots and envelopes remain empty, so the evaluated EU pack still cannot produce legal reassurance.

The contribution is deliberately narrower than production public-key infrastructure. It demonstrates a fail-closed, testable boundary between application code and reviewer-key governance. It does not establish qualified legal review, secure root custody, witnessed publication, hardware-backed monotonicity, or continuity after uninstall or application-data clearing. This distinction continues the thesis's central method: technical evidence must make unsupported authority difficult to imply.


\subsection{Revision 25 contribution: independently witnessed trust-store releases}

Revision 25 prevents a root-signed trust-store envelope from becoming operational solely because its signature verifies. The envelope now stops at \texttt{ENVELOPE\_VERIFIED}; \texttt{CURRENT} additionally requires a separately configured policy with at least two distinct Ed25519 witness anchors and enough valid receipts to meet its threshold. Each receipt binds the exact canonical envelope digest, sequence, witness date, and declared application release identity. Reviewer anchors and rollback state remain unavailable when policy, receipt, or quorum validation fails.

This adds independently attributable cryptographic observations, not a production transparency system. The evaluated root, envelope, witness policy, and receipt sets remain empty. The receipts do not prove witness identity, organisational independence, qualification, uncompromised custody, public inclusion, or the installed APK/AAB bytes. The contribution is therefore a stronger fail-closed release boundary without a corresponding claim of legal authority or deployed assurance.

\input{Iteration-V2/Chapter2/chapter-02-23}

\subsection{Revision 24: key lifecycle and rollback concepts}

Key-management guidance treats cryptographic keys as lifecycle-governed objects rather than timeless constants. NIST SP 800-57 Part 1 Revision 5 discusses general key-management practices and protection across lifecycle phases \cite{nist80057r5}. Privacy Lens uses that source only as security-engineering context: neither NIST guidance nor the project's separation of issuer and approver creates a GDPR requirement or certifies a reviewer.

The Update Framework specifies version, expiry, trusted-root, and rollback-resistance concepts for software-update metadata \cite{tufspec}. Revision 24 borrows the narrower ideas of signed metadata, expiry, monotonic versions, and predecessor continuity, but does not implement TUF roles, threshold signatures, repository metadata, mirrors, consistent snapshots, or a transparency service. Calling the envelope ``TUF compliant'' would therefore exceed the evidence. The academically relevant point is the explicit translation and limitation of external security concepts, not terminological association.


\subsection{Revision 25: attestation subject and digest concepts}

The in-toto Attestation Framework separates an identified subject and digest from a predicate containing claims about that subject \cite{intotoattestation}. Revision 25 borrows the narrower discipline of binding a signed statement to an exact digest and declared release identity. It does not implement the in-toto Statement or envelope formats, a published predicate type, supply-chain layout verification, a transparency log, or in-toto conformance. The receipt-set digest is an order-independent local comparison identifier, not an inclusion proof.

This distinction matters because witness signatures can otherwise invite an unjustified inference that a binary or human role has been independently verified. Privacy Lens receipts currently identify the application by package name, semantic version, version code, and controlled-research channel, but do not include APK, AAB, or signing-certificate digests. Consequently, they can support exact-envelope agreement only; artifact identity and real-world witness governance remain separate acceptance obligations.

\input{Iteration-V2/Chapter3/chapter-03-23}

\subsection{Revision 24 requirements and trust-store threat model}

Revision 24 adds eight requirements. \textbf{TS1} requires a fixed canonical schema whose anchor and revocation ordering cannot alter the signed payload. \textbf{TS2} requires one structurally valid Ed25519 offline root that is valid both at envelope issuance and assessment. \textbf{TS3} forbids an offline-root identifier from also identifying a reviewer anchor and requires different issuer and approver identities. \textbf{TS4} requires explicit, dated revocations within the envelope validity window. \textbf{TS5} requires sequence one to have no predecessor and later sequences to advance exactly one step while binding the accepted predecessor digest. \textbf{TS6} rejects rollback, same-sequence forks, gaps, missing history, invalid signatures, freeze through expiry, and future activation. \textbf{TS7} persists the highest accepted sequence and digest before exposing reviewer anchors and removes all effective anchors if persistence fails. \textbf{TS8} requires the interface to expose both the production provisioning state and local-history limitations.

The attacker model includes mutation of signed fields, substitution of an unknown or malformed root, replay of an older valid envelope, same-sequence replacement, omission of intermediate envelopes, predecessor substitution, and continued use after expiry. The design does not resist an attacker who can replace application code or roots, compromise the device, erase application storage, control the release channel, or subvert the offline ceremony. Those are residual distribution and platform threats, not cases silently counted as passing.


\subsection{Revision 25 witness requirements and attacker model}

Revision 25 adds seven requirements. \textbf{WR1} requires a fixed receipt schema and canonical field order. \textbf{WR2} rejects thresholds below two or above the configured anchor count. \textbf{WR3} requires unique witness owners, key identifiers, and public keys, with identifiers and public-key material distinct from roots and reviewers. \textbf{WR4} requires valid Ed25519 material, coherent validity and revocation dates, and validity at both witnessing and assessment. \textbf{WR5} binds every receipt to the exact envelope digest, sequence, timely witness date, and expected application release identity. \textbf{WR6} rejects the entire supplied set on malformed, duplicate, unknown, revoked, future-dated, wrongly bound, or incorrectly signed receipts. \textbf{WR7} exposes no reviewer anchors and advances no rollback state until quorum passes.

The extended attacker model includes a compromised root acting alone, receipt replay against another envelope or release, duplicate counting, threshold reduction to one, witness/root identity collision, signature mutation, future dating, and use of a revoked witness. It does not claim resistance to colluding or fictitious witnesses, replacement of application code or configured anchors, release-channel compromise, public split views, or device compromise.

\input{Iteration-V2/Chapter4/chapter-04-23}

\subsection{Revision 24 implementation: canonical envelopes and monotonic state}

The implementation serialises a fixed field set, sorts anchors and revocations by key identifier, normalises predecessor digests, and signs the resulting UTF-8 payload with Ed25519. Structural validation checks real calendar dates, safe sequences, exact decoded key and signature lengths, identity separation, duplicate identifiers, validity intervals, and coherent revocation dates. Assessment accepts exactly one matching offline root, checks that it was valid at issuance and remains valid at assessment, verifies the signature, then evaluates stored history before it can expose reviewer anchors.

A repository stores only the highest accepted sequence and SHA-256 envelope digest under a dedicated AsyncStorage key. A pure transition validator rejects decreasing sequences, same-sequence replacement, and skipped steps; the repository fails closed to an invalid assessment with zero anchors when stored state is malformed or cannot be persisted. Clearing ordinary findings does not clear this state. The same boundary is shown in Settings and attached to finding evidence. Because AsyncStorage is app-private rather than tamper-resistant hardware, uninstall, application-data clearing, unsuitable restoration, or device compromise can break continuity. Disabled Android backup narrows one path but does not create a public transparency log.


\subsection{Revision 25 implementation: threshold receipt assessment}

The implementation first performs the Revision-24 envelope assessment. A valid root signature produces \texttt{ENVELOPE\_VERIFIED}, never \texttt{CURRENT}. A second assessor validates the expected package release, witness policy, anchors, and every supplied receipt. Receipt payloads use a fixed JSON field order; SHA-256 identifies the envelope and a canonical witness-key ordering makes the receipt-set digest independent of input order. Ed25519 verifies each canonical receipt. A successful threshold promotes the assessment to \texttt{CURRENT}; every failure path clears effective reviewer anchors and the candidate rollback transition.

The production module intentionally configures no witness policy and an empty receipt set. Settings and expanded finding evidence show the verified/required count and receipt-set digest when available. The interface states that root approval alone is insufficient and that receipts neither bind the installed artifact nor prove witness identity, independence, legal correctness, GDPR compliance, or public transparency. Persistence remains conditional on the combined witnessed assessment.

\input{Iteration-V2/Chapter5/chapter-05-23}

\subsection{Revision 24 evaluation: signed trust-store adversarial probes}

The unchanged seeded compliance corpus returned 800 true positives, 200 true negatives, zero false positives, and zero false negatives. The 1,800-case independent governance oracle also passed. Added probes covered canonical order invariance; signature mutation; unknown, duplicate, malformed, revoked, and not-yet-valid roots; root validity at issuance; envelope expiry and future activation; issuer self-approval; root/reviewer identifier collision; incoherent post-expiry revocation; exact sequence-two advancement; replay of sequence one; a same-sequence fork; predecessor mismatch; sequence gap; missing history; idempotent reassessment; and independent rollback-state transition rejection. Application compilation, accessibility contracts across five files and 12 touchables, release-privacy contracts across 30 source files, and lint passed.

The release build produced a 62,929,891-byte APK with SHA-256 \texttt{\seqsplit{0b723dc0541b93bc27a4ffa5513d1b21f0794d39eaa49a2b719dddd5c7491097}} and a 31,790,208-byte AAB with SHA-256 \texttt{\seqsplit{ccf2e229344941d484db5e582fe8be44940f375641258ae9b7d63d4e74d3f8cf}}. Installation on one OPPO PERM00 reported version code 13, version name 1.12.0, minimum SDK 24, and target SDK 36. Three short force-stop and relaunch rounds returned a live process and the crash buffer remained empty. Device screenshots showed the legal-review warning, the three-link evidence chain, the unprovisioned production trust store, and the local-history boundary without observed clipping in the captured viewport.

These results establish conformance to the stated engineering contract only. The positive cryptographic path uses deterministic synthetic keys; no production root, signed envelope, qualified review, participant outcome, or external replication is claimed.


\subsection{Revision 25 evaluation: witness quorum and release binding}

The fixed-seed compliance corpus again returned 800 true positives, 200 true negatives, zero false positives, and zero false negatives; the 1,800-case independent governance oracle also passed. New adversarial probes covered receipt-set order invariance, missing policy, an invalid one-witness threshold, insufficient quorum, duplicate identities, mutated signatures, a correctly signed wrong envelope digest, a correctly signed wrong release identity, future dating, witness/root collision, current and incoherent revocation, malformed release version, and the invariant that a root-signed envelope alone cannot clear the legal gate. Compilation, accessibility contracts across five files and 12 touchables, release-privacy contracts across 31 source files, and lint passed.

The final release build produced a 62,937,667-byte APK with SHA-256 \texttt{\seqsplit{9e7b6cb5f7fb05a4a9f051be3bbc3e182b509d82bb2d6346480e5ba946513b3a}} and a 31,793,300-byte AAB with SHA-256 \texttt{\seqsplit{bbedcfefa096f1ccf7b75a01a22490cf8f432b2f6af4ea5b38093f5ae682a114}}. Local APK Signature Scheme v2 verification passed. During the preceding build, the OPPO PERM00 rejected a streamed transfer because its received content digest did not verify; non-streaming installation was therefore retained for the final artifact, which succeeded as version code 14 and version name 1.13.0. Three final cold launches took 2,277, 2,266, and 1,937 ms, a process remained live, and the crash buffer was empty.

The captured evidence chain visibly reported production trust store unprovisioned, witnesses 0/2, and production policy unprovisioned. It also rendered the artifact, identity, independence, legal, and transparency limits without observed clipping. These results establish only the stated engineering behavior on one short device sample; deterministic synthetic witness keys do not establish deployed independent witnessing.

\input{Iteration-V2/Chapter6/chapter-06-23}

\subsection{Revision 24 residual validity and governance limits}

Rollback resistance is conditional on retained client state. An uninstall, application-data clearing event, unsuitable restoration, or compromised device can remove or alter the highest accepted sequence and digest. The application also distributes its currently empty roots and envelope through the release itself; it has no independent transparency witness, online status, threshold root control, or recovery protocol. Consequently, the mechanism raises the cost of accidental or local replay inside the admitted execution boundary but does not establish globally consistent trust metadata.

Cryptographic validity remains distinct from legal validity. A valid envelope would establish only that configured key material authorised a particular reviewer-key set. It would not establish the human signer's identity, professional qualification, independence, completeness of review, correctness of source interpretation, or continued applicability to a deployment context. Likewise, the single-device screenshots and three relaunches show a short implemented path, not general usability, comprehension, accessibility, reliability, or OEM coverage. These limitations reserve the main remaining publication-quality claims for external legal, participant, assistive-technology, multi-device, and replication evidence.


\subsection{Revision 25 residual witness and artifact validity limits}

Threshold receipts reduce reliance on one offline root only if witness identities, custody, independence, and ceremonies are genuine. The application cannot infer those organisational facts from public keys. Collusion, coerced approval, shared custody, substituted anchors, or a compromised release process could still produce cryptographically valid receipts. With no public append-only service, gossip, inclusion proof, or split-view detection, different audiences could also receive inconsistent but locally valid evidence.

The current receipt binds a declared release identity rather than artifact digests. It therefore does not prove that the APK/AAB built, distributed, installed, or executed on a device is the release a witness intended to observe. The streamed-install digest failure encountered during evaluation was a transport-path observation, not evidence that the local artifact signature was invalid; successful non-streaming installation likewise does not prove general transport reliability. Production artifact attestation, independent witness governance, qualified legal review, assistive-technology assessment, participant comprehension, multi-device coverage, and external replication remain validity gates.

\input{Iteration-V2/Chapter7/chapter-07-23}

\subsection{Revision 24 product-position update}

Under the strict publication-target rubric introduced after the 70/100 predecessor baseline, Revision 23 was assessed at 94/100 and Revision 24 at 95/100. The one-point increase reflects stronger software and security evidence for reviewer-key governance, not a new legal-validity claim. App 1.12.0 uses version code 13, built as APK and AAB, installed on the recorded OPPO handset, and exposed the unprovisioned production trust store and rollback-state boundary. The score remains an internal expert proxy rather than a university grade, journal decision, legal opinion, accessibility certification, or participant outcome. Production ceremony, qualified review, transparent or witnessed distribution evidence, recovery, and external human studies remain gates rather than awarded credit.

\paragraph{Revision-24 conclusion.} Privacy Lens now validates a canonical Ed25519-signed reviewer trust-store envelope, requires a unique offline root valid at issuance and assessment, represents explicit revocations, and persists a monotonic sequence-and-predecessor chain. Mutation, freeze, rollback, fork, gap, missing-history, malformed-root, and persistence-transition probes fail closed. The evaluated production root list and envelope remain empty, so this additional mechanism cannot convert the unattested EU pack into reassurance. Release and bounded device evidence support the implementation claim, while local-state loss, release-channel control, qualified legal review, human comprehension, accessibility, and external replication remain unresolved. The next legal-priority step is an independently controlled production ceremony with witnessed receipts or transparency evidence and recovery semantics, followed by qualified review and participant evaluation.


\subsection{Revision 25 product-position update}

Under the same strict publication-target rubric, Revision 25 is assessed at 96/100 compared with Revision 24 at 95/100 and the predecessor baseline at 70/100. The one-point increase reflects implemented root-alone resistance, threshold receipt verification, release-identity binding, adversarial coverage, and visible uncertainty. It is not credit for production witnesses, artifact attestation, legal validation, participants, or publication acceptance. A component view is legal compliance 29/30, software and security engineering 25/25, interface and accessibility engineering 20/20, user-psychology safeguards 13/15, and academic evidence 9/10. These are internal expert proxies, not independently validated measurements.

\paragraph{Revision-25 conclusion.} Privacy Lens now requires a valid canonical root-signed envelope and a minimum two-witness receipt threshold before reviewer anchors can become current or rollback state can advance. Receipt replay, duplication, signature mutation, threshold weakening, key-identity collision, revocation, time, envelope, and release-identity mismatches fail closed. The evaluated production trust and witness configuration remains empty, so the mechanism cannot convert the unattested EU pack into reassurance. The next legal-priority step is a genuine independently governed ceremony whose receipts also bind signed release-artifact digests and whose consistency is externally observable, followed by qualified review and independent human and technical evaluation.

{\small
\bibliographystyle{ieeetr}
\bibliography{reference}
}
\end{document}
